% Pulled Pork & Lentil Ragu — Slow Cooker
% Save as: SlowCooker/Pulled_Pork_Lentil_Ragu.tex
% A double batch (8 serves): freeze half.

\begin{recipebox}{Pulled Pork \& Lentil Ragu}{8}{15 min prep + 8--9 hours}

A big batch of pulled pork leg in a smoky tomato sauce, thickened with red
lentils once the pork is shredded. Serve over brown rice, and freeze half for
later.

\begin{ingredients}
  \item[] \textbf{Pork and base}
  \item 1.5\,kg bone-in pork leg, rind and excess fat trimmed
  \item 2 brown onions, diced
  \item 5 garlic cloves, crushed
  \item 3 large carrots, grated
  \item 2 cans (400\,g) chopped tomatoes
  \item 3 tbsp tomato paste
  \item 300\,ml stock (1 chicken or vegetable stock cube in water)
  \item[] \textbf{Seasoning}
  \item 2 tsp dried Italian herbs
  \item 2 tsp smoked paprika
  \item \nicefrac{1}{2} tsp chilli flakes
  \item 1\nicefrac{1}{2} tsp salt
  \item 1 tsp black pepper
  \item[] \textbf{Last 1--1\nicefrac{1}{2} hours}
  \item 200\,g red lentils, rinsed
  \item 250\,ml water
  \item[] \textbf{To finish}
  \item 1 tbsp balsamic vinegar (or any red vinegar)
  \item 150\,g spinach (fresh or frozen)
  \item[] \textbf{To serve}
  \item 400\,g brown rice, uncooked (about 50\,g per serve)
\end{ingredients}

\begin{method}
  \item \textbf{Season the pork.} Mix the smoked paprika, Italian herbs, salt,
        pepper and chilli flakes and rub all over the pork.
  \item \textbf{Build the base.} Put the onions, garlic and grated carrot in the
        slow cooker and sit the pork on top.
  \item \textbf{Add the sauce.} Mix the chopped tomatoes, tomato paste and stock
        and pour around the pork. Start the slow cooker straight away rather
        than leaving the raw pork to sit.
  \item \textbf{Slow cook.} Cook on Low for 7 hours, until the meat pulls away
        from the bone easily.
  \item \textbf{Shred.} Lift the pork onto a board and pull the meat off the bone
        with two forks, discarding the bone, gristle and any fat. Skim any fat
        from the surface of the sauce.
  \item \textbf{Add the lentils.} Stir the lentils and water into the sauce. Cook
        on High with the lid slightly ajar for 1--1\nicefrac{1}{2} hours,
        stirring once halfway, until the lentils are soft and the sauce is thick.
  \item \textbf{Finish.} Return the shredded pork to the pot. Stir in the vinegar
        and spinach until wilted, about 5 minutes. Taste and adjust the salt.
  \item \textbf{Rice.} Cook the brown rice while the pork finishes and serve the
        ragu over it.
\end{method}

\begin{variationbox}{Smoky Chipotle Version}
Swap the Italian herbs for 2 tsp ground cumin and add 1 chipotle in adobo,
finely chopped (or extra chilli flakes). Finish with a squeeze of lime instead
of the vinegar. Serve over rice with salsa, like a Mexican-style bowl.
\end{variationbox}

\begin{variationbox}{Fennel \& Chilli Version}
Replace the Italian herbs with 1 tsp crushed fennel seeds and double the chilli
flakes. This gives a sausage-ragu flavour without the cost of sausages.
\end{variationbox}

\begin{notelist}
  \item A 1.5\,kg bone-in leg may be a tight fit in a smaller slow cooker. If the
        lid won't close, cut the leg in two.
  \item Leg cuts vary. If the pork isn't pulling off the bone at 7 hours, give it
        another hour before shredding.
  \item The lentils go in after the pork is shredded so they don't turn to mush
        over the long cook.
  \item If the sauce is too loose, simmer it on the stovetop for 10 minutes. If
        it's too thick, add a splash of water.
  \item To bulk it out further, add 250\,g chopped mushrooms with the lentils,
        grate in 1--2 courgettes at the end, or stir through frozen peas or mixed
        vegetables in the last 10 minutes.
  \item Also good over mash, in a baked potato or kumara, or as a thick filling
        for wraps and toasties.
  \item Freezes well for up to 3 months. Portion into containers and freeze the
        rice separately.
\end{notelist}

\end{recipebox}
